\subsubsection{ToolOrders}

\begin{figure}[H]
	\centering 
	\includegraphics[width=1\textwidth]{SpecificaTecnica/Tool/ToolOrdineCompleto}
	\caption{Diagramma delle classi per ToolOrders}
\end{figure}

Il diagramma mostra l'architettura completa del tool per la gestione degli ordini.\\
Segue il pattern esagonale: \textit{ToolOrderService} espone la logica applicativa, \textit{ToolOrderAdapter} implementa la porta in ingresso \textit{ToolOrderPortIn}, e \textit{GetOrders} estende \textit{BaseTool} per l'esecuzione effettiva della richiesta all'LLM.\\
Per la definizione di BaseTool, vedere \ref{baseTool}


\subsubsubsection{ToolOrderService}
\begin{figure}[H]
	\centering 
	\includegraphics[width=0.8\textwidth]{SpecificaTecnica/Tool/ToolOrderService}
	\caption{ToolOrders: ToolOrderService}
\end{figure}

\paragraph{Descrizione}
\textit{ToolOrderService} è il service applicativo che espone la logica di recupero dello storico ordini.\\
Viene invocato dall'adapter e si occupa di interrogare il layer di persistenza per restituire una pagina di ordini filtrata per intervallo di date.

\paragraph{Attributi}
\textit{ToolOrderService} non definisce attributi propri.

\paragraph{Metodi e funzioni}
\begin{itemize}
	\item \textbf{\_\_init\_\_(username: string, storico\_service: HistoryService)} - metodo di inizializzazione del service.
	\item \textbf{get\_ordini(pagina: int, data\_inizio: date, data\_fine: date): HistoryPageSchema} - metodo pubblico che recupera una pagina dello storico ordini dell'utente autenticato, filtrata nell'intervallo di date specificato.\\
	Restituisce un oggetto \textit{HistoryPageSchema} contenente la lista degli ordini e le informazioni di paginazione.
\end{itemize}

\subsubsubsection{ToolOrderPortIn}
\begin{figure}[H]
	\centering 
	
\includegraphics[width=0.4\textwidth]{SpecificaTecnica/Tool/ToolOrderPortIn}
	\caption{ToolOrders: ToolOrderPortIn}
\end{figure}

\paragraph{Descrizione}
\textit{ToolOrderPortIn} è l'interfaccia che definisce il contratto per l'accesso allo storico degli ordini da parte dei tool LLM.\\
Viene implementata da \textit{ToolOrderAdapter} e utilizzata da \textit{GetOrders} per disaccoppiare il tool dalla logica applicativa.

\paragraph{Attributi}
\textit{ToolOrderPortIn} non definisce attributi.

\paragraph{Metodi e funzioni}
\begin{itemize}
	\item \textbf{get\_ordini(): HistoryPageSchema} - metodo pubblico che recupera la lista degli ordini.\\
	Viene concretizzato dalla sua implementazione \textit{ToolOrderAdapter}.
\end{itemize}

\subsubsubsection{ToolOrderAdapter}
\begin{figure}[H]
	\centering 
	\includegraphics[width=0.5\textwidth]{SpecificaTecnica/Tool/ToolOrderAdapter}
	\caption{ToolOrders: ToolOrderAdapter}
\end{figure}

\paragraph{Descrizione}
\textit{ToolOrderAdapter} è la classe concreta che implementa l'interfaccia \textit{ToolOrderPortIn}.\\
Funge da adattatore tra il tool LLM$^G$ e il service applicativo, delegando le operazioni effettive a \textit{ToolOrderService}.

\paragraph{Attributi}
\begin{itemize}
	\item \textbf{tool\_order\_service: ToolOrderService} - attributo privato di tipo \textit{ToolOrderService}.
\end{itemize}

\paragraph{Metodi e funzioni}
Oltre alle implementazioni concrete di \textit{ToolOrderPortIn}, è disponibile la seguente funzione:
\begin{itemize}
	\item \textbf{\_\_init\_\_(tool\_order\_service: ToolOrderService)} - metodo di inizializzazione dell'adapter.
\end{itemize}

\subsubsubsection{GetOrders}
\begin{figure}[H]
	\centering 
	\includegraphics[width=0.6\textwidth]{SpecificaTecnica/Tool/GetOrders}
	\caption{ToolOrders: GetOrders}
\end{figure}

\paragraph{Descrizione}
\textit{GetOrders} è il tool concreto che estende \textit{BaseTool} e si occupa di recuperare lo storico degli ordini dell'utente su richiesta del modello LLM.\\
Riceve i parametri di filtro dall'agente, li valida e li delega alla porta \textit{ToolOrderPortIn} per ottenere i dati.

\paragraph{Attributi}
\begin{itemize}
	\item \textbf{tool\_adapter: ToolOrderPortIn} - attributo privato di tipo \textit{ToolOrderPortIn}.\\
	Porta utilizzata per accedere ai dati degli ordini tramite il service applicativo.
\end{itemize}

\paragraph{Metodi e funzioni}
\begin{itemize}
	\item \textbf{\_run(data\_inizio: string, data\_fine: string, pagina: int): string} - metodo pubblico che contiene la logica effettiva di esecuzione: invoca la porta con i parametri ricevuti e formatta la risposta.
	\item \textbf{\_arun(data\_inizio: string, data\_fine: string, pagina: int): string} - versione asincrona di \textit{\_run}.
\end{itemize}

\subsubsubsection{BaseTool}
\begin{figure}[H]
	\centering 
	\includegraphics[width=0.3\textwidth]{SpecificaTecnica/Tool/BaseTool}
	\caption{ToolOrders: BaseTool}
\end{figure}
\label{baseTool}

\paragraph{Descrizione}
\textit{BaseTool} è la classe astratta di base da cui ereditano tutti i tool utilizzati dall'agente LLM.\\
Definisce l'interfaccia comune che ogni tool deve implementare per essere eseguito dall'agente tramite il \textit{ToolExecutor}.

\paragraph{Attributi}
\textit{BaseTool} non definisce attributi propri.

\paragraph{Metodi e funzioni}
\textit{BaseTool} è una classe del framework$^G$ esterno e viene inclusa nel diagramma solo per rappresentare la relazione di ereditarietà. I suoi metodi non sono mostrati perché non rilevanti per la logica del dominio applicativo.


%-----------------------------------------------------------ToolCatalogo ----------------------

\subsubsection{ToolCatalog}
\begin{figure}[H]
	\centering 
	\includegraphics[width=1\textwidth]{SpecificaTecnica/Tool/ToolCatalogoCompleto}
	\caption{ToolCatalog: ToolCatalogoCompleto}
\end{figure}

Il diagramma mostra l'architettura completa del tool per la ricerca nel catalogo prodotti.\\
Per la definizione di BaseTool, vedere sezione \ref{baseTool}

\subsubsubsection{ToolCatalogService}
\label{ToolCatalogService}
\begin{figure}[H]
	\centering
	\includegraphics[width=1\textwidth]{SpecificaTecnica/Tool/ToolCatalogService}
	\caption{Catalogo: ToolCatalogService}
\end{figure}

\paragraph{Descrizione}
Classe che orchestra la ricerca semantica sul catalogo integrando la ricerca
vettoriale con le preferenze di acquisto dell'utente.\\
Recupera i prodotti tramite \textit{EmbeddedCatalogService} e li riordina privilegiando
quelli acquistati con maggiore frequenza dall'utente, qualora siano disponibili le
preferenze.

\paragraph{Attributi}
\begin{itemize}
	\item \textbf{embedded\_catalog: EmbeddedCatalogService} - istanza del servizio di ricerca vettoriale sul catalogo.\\
	Per la definizione di EmbeddedCatalogService, vedere \ref{EmbeddedCartService}.
	\item \textbf{catalog\_repo: CatalogRepoPort} - istanza del porto di accesso ai dati del catalogo, utilizzata per recuperare i dettagli dei prodotti trovati.
	Per la definizione di CatalogRepoPort, vedere \ref{CatalogRepoPort}.
	\item \textbf{preferred\_product\_frequency: Dict[string, int]} - dizionario che
	mappa gli identificativi di prodotto alla frequenza di acquisto dell'utente
	corrente; valore di default è un dizionario vuoto.
\end{itemize}

\paragraph{Metodi e funzioni}
\begin{itemize}
	\item \textbf{\_\_init\_\_(embedded\_catalog: EmbeddedCatalogService, catalog\_repo: CatalogRepoPort, preferred\_product\_frequency: Dict[string, int])} - inizializza il servizio ricevendo le dipendenze; se \textit{preferred\_product\_frequency} è \textit{None}, viene sostituito con un dizionario vuoto.\\
	Per la definizione di CatalogRepoPort, vedere \ref{CatalogRepoPort}.
	\item \textbf{search\_catalog(query: string, threshold: float): List[CatalogProduct]}
	- ricerca i prodotti semanticamente simili alla query tramite
	\textit{embedded\_catalog}, ne recupera i dettagli tramite
	\textit{catalog\_repo} e, se sono disponibili le preferenze utente, riordina i
	risultati privilegiando i prodotti con frequenza di acquisto maggiore, mantenendo
	l'ordine di similarità come criterio secondario.
\end{itemize}

\subsubsubsection{ToolCatalogPortIn}
\begin{figure}[H]
	\centering 
	
\includegraphics[width=0.7\textwidth]{SpecificaTecnica/Tool/ToolCatalogPortIn}
	\caption{ToolCatalog: ToolCatalogPortIn}
\end{figure}

\paragraph{Descrizione}
\textit{ToolCatalogPortIn} è l'interfaccia che definisce il contratto per la ricerca nel catalogo prodotti da parte dei tool LLM.\\
Viene implementata da \textit{ToolCatalogAdapter} e utilizzata da \textit{SearchCatalog} per disaccoppiare il tool dalla logica applicativa.

\paragraph{Attributi}
\textit{ToolCatalogPortIn} non definisce attributi.

\paragraph{Metodi e funzioni}
\begin{itemize}
	\item \textbf{search\_catalog(query: string, threshold: float): List[CatalogProduct]} - metodo pubblico che effettua la ricerca nel catalogo.\\
	Viene concretizzato dalla sua implementazione \textit{ToolCatalogAdapter}.
\end{itemize}

\subsubsubsection{ToolCatalogAdapter}
\begin{figure}[H]
	\centering 
	\includegraphics[width=0.6\textwidth]{SpecificaTecnica/Tool/ToolCatalogAdapter}
	\caption{ToolCatalog: ToolCatalogAdapter}
\end{figure}

\paragraph{Descrizione}
\textit{ToolCatalogAdapter} è la classe concreta che implementa l'interfaccia \textit{ToolCatalogPortIn}.\\
Funge da adattatore tra il tool LLM$^G$ e il service applicativo, delegando le operazioni di ricerca a \textit{ToolCatalogService}.

\paragraph{Attributi}
\begin{itemize}
	\item \textbf{tool\_catalog\_service: ToolCatalogService} - attributo privato di tipo \textit{ToolCatalogService}.
\end{itemize}

\paragraph{Metodi e funzioni}
Oltre alle implementazioni concrete di \textit{ToolCatalogPortIn}, è disponibile la seguente funzione:
\begin{itemize}
	\item \textbf{\_\_init\_\_(tool\_catalog\_service: ToolCatalogService)} - metodo di inizializzazione dell'adapter.
\end{itemize}

\subsubsubsection{SearchCatalogTool}
\begin{figure}[H]
	\centering 
	\includegraphics[width=0.4\textwidth]{SpecificaTecnica/Tool/SearchCatalog}
	\caption{ToolCatalog: SearchCatalog}
\end{figure}

\paragraph{Descrizione}
\textit{SearchCatalog} è il tool concreto che estende \textit{BaseTool} e consente all'agente LLM$^G$ di ricercare prodotti nel catalogo.\\
Riceve la query$^G$ dall'agente, la delega alla porta \textit{ToolCatalogPortIn} e restituisce i risultati serializzati.

\paragraph{Attributi}
\begin{itemize}
	\item \textbf{tool\_adapter: ToolCatalogPortIn} - attributo privato di tipo \textit{ToolCatalogPortIn}.\\
	Porta utilizzata per accedere alla logica di ricerca nel catalogo tramite il service applicativo.
\end{itemize}

\paragraph{Metodi e funzioni}
\begin{itemize}
	\item \textbf{\_run(query$^G$: string, threshold: float): string} - metodo pubblico che contiene la logica effettiva di esecuzione: invoca la porta con la query$^G$ ricevuta e formatta la risposta.
	\item \textbf{\_arun(query$^G$: string, threshold: float): string} - versione asincrona di \textit{\_run}.
\end{itemize}


%-----------------------------------------------------------Tool Cart ----------------------
\subsubsection{ToolCart}
\begin{figure}[H]
	\centering 
	\includegraphics[width=1\textwidth]{SpecificaTecnica/Tool/ToolCartCompleto}
	\caption{ToolCart: ToolCartCompleto}
\end{figure}

Il diagramma mostra l'architettura completa del tool per la gestione del carrello.\\
Per la definizione di BaseTool, vedere sezione \ref{baseTool}

\subsubsubsection{ToolCartService}
\begin{figure}[H]
	\centering 
	\includegraphics[width=0.9\textwidth]{SpecificaTecnica/Tool/ToolCartService}
	\caption{ToolCart: ToolCartService}
\end{figure}
\label{ToolCartService}

\paragraph{Descrizione}
\textit{ToolCartService} è il service applicativo che espone tutte le operazioni disponibili sul carrello dell'utente.\\
Viene invocato dall'adapter e si occupa di delegare le operazioni al layer di persistenza.

\paragraph{Attributi}
\textit{ToolCartService} non definisce attributi propri.

\paragraph{Metodi e funzioni}
\begin{itemize}
	\item \textbf{\_\_init\_\_(username: string, cart\_repo: CartRepoPort, embedded\_cart: EmbeddedCartService)} - metodo di inizializzazione del service. Per la definizione di CartRepoPort vedere \ref{CartRepoPort}.\\
	Per la definizione di EmbeddedCartService, vedere \ref{EmbeddedCartService}
	\item \textbf{get\_cart\_items(): List[CartProduct]} - metodo pubblico che recupera la lista di tutti i prodotti presenti nel carrello dell'utente corrente.
	\item \textbf{search\_cart(query: string, threshold: float): List[CartProduct]} - metodo pubblico che effettua una ricerca semantica tra i prodotti nel carrello in base alla query e alla soglia di similarità specificate.\\
	Restituisce la lista dei prodotti che superano la soglia.
	\item \textbf{add\_to\_cart(prod\_id: string, qty: int): CartProduct} - metodo pubblico che aggiunge al carrello il prodotto identificato da \textit{prod\_id} con la quantità specificata.\\
	Restituisce il prodotto aggiunto al carrello.
	\item \textbf{remove\_from\_cart(prod\_id: string): CartProduct} - metodo pubblico che rimuove dal carrello il prodotto identificato da \textit{prod\_id}.\\
	Restituisce il prodotto rimosso.
	\item \textbf{update\_cart\_qty(prod\_id: string, qty: int, operation: UpdateCartOperation): CartProduct} - metodo pubblico che aggiorna la quantità del prodotto identificato da \textit{prod\_id} nel carrello, applicando l'operazione specificata.\\
	Restituisce il prodotto aggiornato.
\end{itemize}

\subsubsubsection{ToolCartPortIn}
\begin{figure}[H]
	\centering 
	
\includegraphics[width=0.9\textwidth]{SpecificaTecnica/Tool/ToolCartPortIn}
	\caption{ToolCart: ToolCartPortIn}
\end{figure}

\paragraph{Descrizione}
\textit{ToolCartPortIn} è l'interfaccia che definisce il contratto per tutte le operazioni sul carrello da parte dei tool LLM.\\
Viene implementata da \textit{ToolCartAdapter} e utilizzata dai tool concreti per disaccoppiare la logica di esecuzione dal layer applicativo.

\paragraph{Attributi}
\textit{ToolCartPortIn} non definisce attributi.

\paragraph{Metodi e funzioni}
\begin{itemize}
	\item \textbf{add\_to\_cart(prod\_id: string, qty: int): string} - metodo pubblico che aggiunge un prodotto al carrello.
	\item \textbf{remove\_from\_cart(prod\_id: string): string} - metodo pubblico che rimuove un prodotto dal carrello.
	\item \textbf{update\_cart\_item\_qty(prod\_id: string, qty: int, operation: CartUpdateOperation): string} - metodo pubblico che aggiorna la quantità di un prodotto nel carrello.
	\item \textbf{get\_cart\_items(): List[CartProduct]} - metodo pubblico che recupera la lista dei prodotti nel carrello.
	\item \textbf{search\_cart(query: string, threshold: float): List[CartProduct]} - metodo pubblico che effettua una ricerca semantica tra i prodotti nel carrello.
\end{itemize}

\subsubsubsection{ToolCartAdapter}
\begin{figure}[H]
	\centering 
	\includegraphics[width=0.6\textwidth]{SpecificaTecnica/Tool/ToolCartAdapter}
	\caption{ToolCart: ToolCartAdapter}
\end{figure}

\paragraph{Descrizione}
\textit{ToolCartAdapter} è la classe concreta che implementa l'interfaccia \textit{ToolCartPortIn}.\\
Funge da adattatore tra i tool LLM$^G$ e il service applicativo, delegando tutte le operazioni sul carrello a \textit{ToolCartService}.

\paragraph{Attributi}
\begin{itemize}
	\item \textbf{tool\_cart\_service: ToolCartService} - attributo privato di tipo \textit{ToolCartService}.
\end{itemize}

\paragraph{Metodi e funzioni}
Oltre alle implementazioni concrete di \textit{ToolCartPortIn}, è disponibile la seguente funzione:
\begin{itemize}
	\item \textbf{\_\_init\_\_(tool\_cart\_service: ToolCartService)} - metodo di inizializzazione dell'adapter.
\end{itemize}

\subsubsubsection{AddToCart}
\begin{figure}[H]
	\centering 
	\includegraphics[width=0.7\textwidth]{SpecificaTecnica/Tool/AddToCart}
	\caption{ToolCart: AddToCart}
\end{figure}

\paragraph{Descrizione}
\textit{AddToCart} è il tool concreto che estende \textit{BaseTool} e consente all'agente LLM$^G$ di aggiungere un prodotto al carrello dell'utente.\\
Riceve l'identificativo del prodotto e la quantità dall'agente e delega l'operazione alla porta \textit{ToolCartPortIn}.

\paragraph{Attributi}
\begin{itemize}
	\item \textbf{tool\_adapter: ToolCartPortIn} - attributo privato di tipo \textit{ToolCartPortIn}.\\
	Porta utilizzata per accedere alle operazioni sul carrello tramite il service applicativo.
\end{itemize}

\paragraph{Metodi e funzioni}
\begin{itemize}
	\item \textbf{\_run(prod\_id: string, qty: int): string} - metodo pubblico che contiene la logica effettiva di esecuzione: invoca la porta per aggiungere il prodotto specificato al carrello dell'utente.
	\item \textbf{\_arun(prod\_id: string, qty: int): string} - versione asincrona di \textit{\_run}.
\end{itemize}

\subsubsubsection{RemoveFromCart}
\begin{figure}[H]
	\centering 
	\includegraphics[width=0.5\textwidth]{SpecificaTecnica/Tool/RemoveFromCart}
	\caption{ToolCart: RemoveFromCart}
\end{figure}

\paragraph{Descrizione}
\textit{RemoveFromCart} è il tool concreto che estende \textit{BaseTool} e consente all'agente LLM$^G$ di rimuovere un prodotto dal carrello dell'utente.\\
Riceve l'identificativo del prodotto dall'agente e delega la rimozione alla porta \textit{ToolCartPortIn}.

\paragraph{Attributi}
\begin{itemize}
	\item \textbf{tool\_adapter: ToolCartPortIn} - attributo privato di tipo \textit{ToolCartPortIn}.\\
	Porta utilizzata per accedere alle operazioni sul carrello tramite il service applicativo.
\end{itemize}

\paragraph{Metodi e funzioni}
\begin{itemize}
	\item \textbf{\_run(prod\_id: string): string} - metodo pubblico che contiene la logica effettiva di esecuzione: invoca la porta per rimuovere il prodotto specificato dal carrello dell'utente.
	\item \textbf{\_arun(prod\_id: string): string} - versione asincrona di \textit{\_run}.
\end{itemize}

\subsubsubsection{UpdateCartItemQty}
\begin{figure}[H]
	\centering 
	\includegraphics[width=0.7\textwidth]{SpecificaTecnica/Tool/UpdateCartItemQty}
	\caption{ToolCart: UpdateCartItemQty}
\end{figure}

\paragraph{Descrizione}
\textit{UpdateCartItemQty} è il tool concreto che estende \textit{BaseTool} e consente all'agente LLM$^G$ di aggiornare la quantità di un prodotto presente nel carrello dell'utente.\\
Riceve l'identificativo del prodotto, la nuova quantità e il tipo di operazione dall'agente, delegando l'aggiornamento alla porta \textit{ToolCartPortIn}.

\paragraph{Attributi}
\begin{itemize}
	\item \textbf{tool\_adapter: ToolCartPortIn} - attributo privato di tipo \textit{ToolCartPortIn}.\\
	Porta utilizzata per accedere alle operazioni sul carrello tramite il service applicativo.
\end{itemize}

\paragraph{Metodi e funzioni}
\begin{itemize}
	\item \textbf{\_run(prod\_id: string, qty: int, operation: CartUpdateOperation): string} - metodo pubblico che contiene la logica effettiva di esecuzione: invoca la porta per aggiornare la quantità del prodotto specificato nel carrello dell'utente.
	\item \textbf{\_arun(prod\_id: string, qty: int, operation: CartUpdateOperation): string} - versione asincrona di \textit{\_run}.
\end{itemize}

\subsubsubsection{SearchCart}
\begin{figure}[H]
	\centering 
	\includegraphics[width=0.6\textwidth]{SpecificaTecnica/Tool/SearchCart}
	\caption{ToolCart: SearchCart}
\end{figure}

\paragraph{Descrizione}
\textit{SearchCart} è il tool concreto che estende \textit{BaseTool} e consente all'agente LLM$^G$ di effettuare una ricerca semantica tra i prodotti presenti nel carrello dell'utente.\\
Riceve la query$^G$ dall'agente e delega la ricerca alla porta \textit{ToolCartPortIn}.

\paragraph{Attributi}
\begin{itemize}
	\item \textbf{tool\_adapter: ToolCartPortIn} - attributo privato di tipo \textit{ToolCartPortIn}.\\
	Porta utilizzata per accedere alle operazioni sul carrello tramite il service applicativo.
\end{itemize}

\paragraph{Metodi e funzioni}
\begin{itemize}
	\item \textbf{\_run(query$^G$: string, threshold: float): string} - metodo pubblico che contiene la logica effettiva di esecuzione: invoca la porta con la query$^G$ ricevuta e formatta la risposta.
	\item \textbf{\_arun(query$^G$: string, threshold: float): string} - versione asincrona di \textit{\_run}.
\end{itemize}

%-----------------------------------------------------------LLMAgent ----------------------
\subsubsection{LLM Agent}
\begin{figure}[H]
	\centering 
	\includegraphics[width=0.5\textwidth]{SpecificaTecnica/Tool/AgentCompleto}
	\caption{Diagramma delle classi per LLMAgent}
\end{figure}

Il diagramma mostra l'architettura completa dell'agente LLM.\\
Per la definizione di BaseTool, vedere sezione \ref{baseTool}\\
Per la definizione di ChatService, vedere sezione \ref{ChatService}

\subsubsubsection{ToolExecutor}
\begin{figure}[H]
	\centering 
	\includegraphics[width=0.5\textwidth]{SpecificaTecnica/Tool/ToolExecutor}
	\caption{LLM Agent: ToolExecutor}
\end{figure}
\label{ToolExecutor}

\paragraph{Descrizione}
\textit{ToolExecutor} è la classe responsabile$^G$ della selezione e dell'esecuzione dei tool disponibili per l'agente LLM.\\
Mantiene internamente la lista di tutti i \textit{BaseTool} registrati e, dato il nome di un tool e i relativi argomenti, ne esegue il metodo \textit{run}.

\paragraph{Attributi}
\begin{itemize}
	\item \textbf{tools: List[BaseTool]} - attributo privato che rappresenta la lista di tutti i tool disponibili per l'agente.
\end{itemize}

\paragraph{Metodi e funzioni}
\begin{itemize}
	\item \textbf{\_\_init\_\_(tools: List[BaseTool])} - metodo di inizializzazione dell'executor.
	\item \textbf{execute(tool\_name: string): string} - metodo pubblico che individua il tool corrispondente al nome fornito e ne invoca il metodo \textit{run} con gli argomenti specificati.\\
	Restituisce il risultato dell'esecuzione come stringa.
	\item \textbf{get\_tool\_by\_name(tool\_name: string): BaseTool} - metodo pubblico che restituisce il \textit{BaseTool} registrato corrispondente al nome fornito.
\end{itemize}

\subsubsubsection{LLMAgent}
\begin{figure}[H]
	\centering 
	\includegraphics[width=0.7\textwidth]{SpecificaTecnica/Tool/LLMAgent}
	\caption{LLM Agent: LLMAgent}
\end{figure}
\label{LLMAgent}

\paragraph{Descrizione}
\textit{LLMAgent} è il componente centrale dell'architettura agente.\\
Si occupa di gestire l'intero ciclo di elaborazione: costruisce il blocco di preferenze utente, esegue il loop agentico interagendo con il modello LLM$^G$ e i tool disponibili tramite \textit{ToolExecutor}, normalizza il contenuto della risposta e produce un oggetto \textit{LLMResponse} da restituire al chiamante.

\paragraph{Attributi}
\begin{itemize}
	\item \textbf{model: ChatOpenAI} - attributo privato di tipo \textit{ChatOpenAI}.\\
	Rappresenta il client del modello di linguaggio utilizzato per generare le risposte.
	\item \textbf{toolExecutor: ToolExecutor} - attributo privato di tipo \textit{ToolExecutor}.\\
	Componente utilizzato per selezionare ed eseguire i tool richiesti dal modello durante il loop agentico.
\end{itemize}

\paragraph{Metodi e funzioni}
\begin{itemize}
	\item \textbf{\_\_init\_\_(tool\_executor: ToolExecutor)} - metodo di inizializzazione dell'agente.
	\item \textbf{\_build\_user\_pref\_block(preferences: List[UserPreferences]): string} - metodo privato che costruisce il blocco testuale delle preferenze utente da includere nel prompt di sistema.
	\item \textbf{\_run\_agentic\_loop(messages: List[BaseMessage]): AgenticLoopResult} - metodo privato che esegue il loop agentico: invia i messaggi al modello, intercetta le richieste di tool call, le esegue tramite \textit{ToolExecutor} e ripete finché il modello non produce una risposta finale.
	\item \textbf{\_build\_response(result: AgenticLoopResult): LLMResponse} - metodo privato che trasforma il risultato del loop agentico in un oggetto \textit{LLMResponse} strutturato.
	\item \textbf{\_normalize\_content(content: NormalizableContent): string} - metodo privato che normalizza il contenuto della risposta del modello, convertendolo in una stringa uniforme indipendentemente dal tipo originale.
	\item \textbf{\_build\_history(request: LLMRequest): List[BaseMessage]} - metodo privato che costruisce la cronologia dei messaggi da inviare al modello a partire dalla richiesta ricevuta.
	\item \textbf{\_clean(text: string): string} - metodo privato che effettua operazioni di pulizia sul testo prodotto dal modello, rimuovendo artefatti o formattazioni indesiderate.
	\item \textbf{invoke(request: LLMRequest): LLMResponse} - metodo pubblico che avvia l'elaborazione della richiesta: costruisce la cronologia, esegue il loop agentico e restituisce la risposta finale.
\end{itemize}

\subsubsubsection{LLMPort}
\begin{figure}[H]
	\centering 
	
\includegraphics[width=0.6\textwidth]{SpecificaTecnica/Tool/LLMPort}
	\caption{LLM Agent: LLMPort}
\end{figure}
\label{LLMPort}

\paragraph{Descrizione}
\textit{LLMPort} è l'interfaccia che definisce il contratto per l'invocazione dell'agente LLM$^G$ da parte del layer applicativo.\\
Viene implementata da \textit{LLMAdapter} e utilizzata da \textit{ChatService} per disaccoppiare la logica di business dall'implementazione concreta dell'agente.

\paragraph{Attributi}
\textit{LLMPort} non definisce attributi.

\paragraph{Metodi e funzioni}
\begin{itemize}
	\item \textbf{invoke(request: LLMRequest): LLMResponse} - metodo pubblico che riceve la richiesta e restituisce la risposta elaborata dall'agente LLM.\\
	Viene concretizzato dalla sua implementazione \textit{LLMAdapter}.
\end{itemize}

\subsubsubsection{LLMAdapter}
\begin{figure}[H]
	\centering 
	\includegraphics[width=0.6\textwidth]{SpecificaTecnica/Tool/LLMAdapter}
	\caption{LLM Agent: LLMAdapter}
\end{figure}
\label{LLMAdapter}

\paragraph{Descrizione}
\textit{LLMAdapter} è la classe concreta che implementa l'interfaccia \textit{LLMPort}.\\
Funge da adattatore tra il layer applicativo e l'agente LLM$^G$, delegando l'elaborazione delle richieste all'istanza di \textit{LLMAgent} che incapsula al proprio interno.

\paragraph{Attributi}
\begin{itemize}
	\item \textbf{agent: LLMAgent} - attributo privato di tipo \textit{LLMAgent}.\\
	Rappresenta l'agente LLM$^G$ utilizzato per elaborare le richieste ricevute.
\end{itemize}

\paragraph{Metodi e funzioni}
Oltre alle implementazioni concrete di \textit{LLMPort}, è disponibile la seguente funzione:
\begin{itemize}
	\item \textbf{\_\_init\_\_(agent: LLMAgent)} - metodo di inizializzazione dell'adapter.
\end{itemize}